% =====================================================================
%  PREÁMBULO
% 
% Adiciones y modificaciones a 'ucmthesis.cls'
% =====================================================================

% ----------------------- Codificación e idioma ----------------------
\usepackage[utf8]{inputenc}
\usepackage[T1]{fontenc}
\IfFileExists{lmodern.sty}{\usepackage{lmodern}}{} % fuente Latin Modern (si está disponible)
\usepackage[spanish,es-noshorthands]{babel}
\addto\captionsspanish{%
  \renewcommand{\contentsname}{Tabla de contenidos}%
  }

% ----------------------- Matemáticas ---------------------------------
\usepackage{amsmath,amssymb,amsthm}

% ----------------------- Gráficos y tablas ---------------------------
\usepackage{graphicx}
\usepackage{xcolor}
\usepackage{booktabs}
\usepackage{array}
\usepackage{tikz}
\usetikzlibrary{arrows.meta,positioning,calc}

% ----------------------- Código Lean (listings) ----------------------
\usepackage{listings}

\usepackage[strings]{underscore}

\definecolor{leankw}{RGB}{0,90,180}       
\definecolor{leantactic}{RGB}{130,30,120} 
\definecolor{leancomment}{RGB}{110,110,110}
\definecolor{leanbg}{RGB}{247,247,250}
\definecolor{leanframe}{RGB}{200,200,210}

\lstdefinelanguage{lean}{
  keywords={import, inductive, structure, where, def, theorem, example,
            local, infixr, infixl, prefix, variable, extends, fun, match,
            by, with, deriving, open, namespace, end, noncomputable},
  keywords=[2]{intro, intros, exact, apply, have, constructor, rfl,
            by_contra, by_cases, contradiction, rcases, rintro, obtain,
            refine, induction, subst, simp, generalizing, at, cases,
            absurd, trivial, assumption},
  sensitive=true,
  comment=[l]{--},
  morecomment=[s]{/-}{-/},
  literate=
    {→}{{$\to$}}1               {↔}{{$\leftrightarrow$}}1
    {⟶}{{$\longrightarrow$}}2  {⟷}{{$\longleftrightarrow$}}2
    {⟹}{{$\Longrightarrow$}}2  {⟸}{{$\Longleftarrow$}}2
    {∼}{{$\sim$}}1              {¬}{{$\lnot$}}1
    {⋏}{{$\curlywedge$}}1       {⋎}{{$\curlyvee$}}1
    {∧}{{$\land$}}1             {∨}{{$\lor$}}1
    {□}{{$\square$}}1           {◇}{{$\lozenge$}}1
    {⊥}{{$\bot$}}1              {⊤}{{$\top$}}1
    {⊢}{{$\vdash$}}1            {⊬}{{$\nvdash$}}1
    {⊨}{{$\vDash$}}1
    {∈}{{$\in$}}1               {∉}{{$\notin$}}1
    {⊆}{{$\subseteq$}}1         {∪}{{$\cup$}}1
    {⋃}{{$\bigcup$}}1           {∅}{{$\varnothing$}}1
    {∀}{{$\forall$}}1           {∃}{{$\exists$}}1
    {⟨}{{$\langle$}}1            {⟩}{{$\rangle$}}1
    {·}{{$\cdot$}}1
    {φ}{{$\varphi$}}1           {ψ}{{$\psi$}}1
    {χ}{{$\chi$}}1              {ξ}{{$\xi$}}1
    {Γ}{{$\Gamma$}}1            {Δ}{{$\Delta$}}1
    {Θ}{{$\Theta$}}1
    {α}{{$\alpha$}}1              {β}{{$\beta$}}1
    {₀}{{${}_0$}}1              {₁}{{${}_1$}}1
    {₂}{{${}_2$}}1              {ₙ}{{${}_n$}}1
    {á}{{\'a}}1 {é}{{\'e}}1 {í}{{\'{\i}}}1 {ó}{{\'o}}1 {ú}{{\'u}}1
    {ñ}{{\~n}}1 {ü}{{\"u}}1
    {Á}{{\'A}}1 {Í}{{\'I}}1 {Ó}{{\'O}}1 {Ú}{{\'U}}1,
}

\lstset{
  language=lean,
  basicstyle=\ttfamily\small,
  keywordstyle=\color{leankw}\bfseries,
  keywordstyle=[2]\color{leantactic},
  commentstyle=\color{leancomment}\itshape,
  backgroundcolor=\color{leanbg},
  frame=single,
  rulecolor=\color{leanframe},
  framesep=5pt,
  xleftmargin=8pt, xrightmargin=8pt,
  showstringspaces=false,
  breaklines=true,
  breakatwhitespace=true,
  columns=fullflexible,
  keepspaces=true,
  aboveskip=0.9em, belowskip=0.9em,
    numbers=left,
  numberstyle=\tiny\color{leancomment},
  numbersep=8pt,
  firstnumber=1,
}

% --- Símbolos Unicode utilizables en modo texto (para \lean y la prosa) ---
\DeclareUnicodeCharacter{2192}{\ensuremath{\to}}
\DeclareUnicodeCharacter{2194}{\ensuremath{\leftrightarrow}}
\DeclareUnicodeCharacter{27F6}{\ensuremath{\longrightarrow}}
\DeclareUnicodeCharacter{27F7}{\ensuremath{\longleftrightarrow}}
\DeclareUnicodeCharacter{27F9}{\ensuremath{\Longrightarrow}}
\DeclareUnicodeCharacter{27F8}{\ensuremath{\Longleftarrow}}
\DeclareUnicodeCharacter{223C}{\ensuremath{\sim}}
\DeclareUnicodeCharacter{00AC}{\ensuremath{\lnot}}
\DeclareUnicodeCharacter{22CF}{\ensuremath{\curlywedge}}
\DeclareUnicodeCharacter{22CE}{\ensuremath{\curlyvee}}
\DeclareUnicodeCharacter{2227}{\ensuremath{\land}}
\DeclareUnicodeCharacter{2228}{\ensuremath{\lor}}
\DeclareUnicodeCharacter{25A1}{\ensuremath{\square}}
\DeclareUnicodeCharacter{25C7}{\ensuremath{\lozenge}}
\DeclareUnicodeCharacter{22A5}{\ensuremath{\bot}}
\DeclareUnicodeCharacter{22A4}{\ensuremath{\top}}
\DeclareUnicodeCharacter{22A2}{\ensuremath{\vdash}}
\DeclareUnicodeCharacter{22AC}{\ensuremath{\nvdash}}
\DeclareUnicodeCharacter{22A8}{\ensuremath{\vDash}}
\DeclareUnicodeCharacter{2208}{\ensuremath{\in}}
\DeclareUnicodeCharacter{2209}{\ensuremath{\notin}}
\DeclareUnicodeCharacter{2286}{\ensuremath{\subseteq}}
\DeclareUnicodeCharacter{222A}{\ensuremath{\cup}}
\DeclareUnicodeCharacter{22C3}{\ensuremath{\bigcup}}
\DeclareUnicodeCharacter{2205}{\ensuremath{\varnothing}}
\DeclareUnicodeCharacter{2200}{\ensuremath{\forall}}
\DeclareUnicodeCharacter{2203}{\ensuremath{\exists}}
\DeclareUnicodeCharacter{27E8}{\ensuremath{\langle}}
\DeclareUnicodeCharacter{27E9}{\ensuremath{\rangle}}
\DeclareUnicodeCharacter{00B7}{\ensuremath{\cdot}}
\DeclareUnicodeCharacter{03C6}{\ensuremath{\varphi}}
\DeclareUnicodeCharacter{03C8}{\ensuremath{\psi}}
\DeclareUnicodeCharacter{03C7}{\ensuremath{\chi}}
\DeclareUnicodeCharacter{03BE}{\ensuremath{\xi}}
\DeclareUnicodeCharacter{0393}{\ensuremath{\Gamma}}
\DeclareUnicodeCharacter{0394}{\ensuremath{\Delta}}
\DeclareUnicodeCharacter{0398}{\ensuremath{\Theta}}
\DeclareUnicodeCharacter{03B1}{\ensuremath{\alpha}}
\DeclareUnicodeCharacter{03B2}{\ensuremath{\beta}}
\DeclareUnicodeCharacter{2080}{\ensuremath{{}_0}}
\DeclareUnicodeCharacter{2081}{\ensuremath{{}_1}}
\DeclareUnicodeCharacter{2082}{\ensuremath{{}_2}}
\DeclareUnicodeCharacter{2099}{\ensuremath{{}_n}}

% Código Lean en línea: \lean{...}
\newcommand{\lean}[1]{\texttt{\small #1}}

% ----------------------- Hipervínculos -------------------------------
\usepackage[colorlinks=true, linkcolor=black, citecolor=black,
            urlcolor=blue!60!black,
            bookmarksnumbered=true,
            pdftitle={Lógica modal en Lean},
            pdfauthor={Jorge Martinena Cepa}]{hyperref}

% ----------------------- Entornos de teoremas ------------------------
\theoremstyle{plain}
\newtheorem{teorema}{Teorema}[chapter]
\newtheorem{lema}[teorema]{Lema}
\newtheorem{proposicion}[teorema]{Proposición}
\newtheorem{corolario}[teorema]{Corolario}

\theoremstyle{definition}
\newtheorem{definicion}[teorema]{Definición}
\newtheorem{ejemplo}[teorema]{Ejemplo}

\theoremstyle{remark}
\newtheorem{observacion}[teorema]{Observación}

% ----------------------- Macros de lógica modal ----------------------
\newcommand{\nec}{\square}          % operador de necesidad
\newcommand{\pos}{\lozenge}         % operador de posibilidad
\newcommand{\limp}{\to}             % implicación (objeto)
\newcommand{\proves}{\vdash}        % demostrabilidad
\newcommand{\sat}{\vDash}           % satisfacción M,w |=
\newcommand{\K}{\mathbf{K}}         % el sistema K

% ----------------------- Corrección manual de la portada --------------
\makeatletter
\renewcommand{\maketitle}{%
\begin{alwayssingle}
    \renewcommand{\footnotesize}{\small}
    \renewcommand{\footnoterule}{\relax}
    \thispagestyle{empty}
  \null\vfill
  \begin{center}
  {\noindent\rule{15cm}{1.5pt}\par}
  \vspace*{2.4mm}
    { \Huge {\bfseries {\@title}} \par}
  \vspace*{3.6mm}
  {\noindent\rule{15cm}{1.5pt}\par}
{\large \vspace*{18mm} {\logo \par} \vspace*{22.5mm}}
    {{\Large \@author} \par}
{\large \vspace*{1ex}
    {{\@college} \par}
\vspace*{1ex}
    {Universidad Complutense de Madrid \par}
\vspace*{25mm}
    {{\submittedtext} \par}
\vspace*{1ex}
    {\it {\@degree} \par}
\vspace*{2ex}
    {\@degreedate}}
  \end{center}
  \null\vfill
\end{alwayssingle}}
\makeatother

